\solchapter{accel}

\sol{accel}{1} $p = h/\lambda$; with $\lambda = 0.1\un{fm}$, $pc = 2\pi\hbar c/\lambda = 12\un{GeV}$ (or $2\un{GeV}$ using $\hbar/\lambda$; the factor of $2\pi$ depends on what ``comparable'' means). At a few TeV per parton, $\hbar c/E\approx10^{-4}\un{fm} = 10^{-19}\un{m}$: ten thousand times smaller than a proton.

\sol{accel}{2} $f = eB/2\pi m_p = 23\un{MHz}$. At the edge, $p = eBR$, so $K = (eBR)^2/2m_p = 27\un{MeV}$. A few percent of $m_pc^2$ is $20$--$30\un{MeV}$, so this cyclotron is right at the limit. The proton's frequency drops as $1/\gamma$ and the particles fall out of step; synchrocyclotrons fix this by lowering the frequency as the protons speed up.

\sol{accel}{3} $B = p/0.3\rho = 6800/(0.3\times2800) = 8.1\un{T}$. A $100\un{km}$ ring with the same fill factor has $\rho = \tfrac23(100\un{km})/2\pi = 10.6\un{km}$, and $50\un{TeV}$ needs $B = 50\,000/(0.3\times10\,600) = 16\un{T}$ --- the design goal of the proposed Future Circular Collider.

\sol{accel}{4} The invariant is $s = (E + mc^2)^2 - p^2c^2$, which equals $2Emc^2 + 2m^2c^4\approx2Emc^2$. Fixed target: $\sqrt{2(6800)(0.938)} = 113\un{GeV}$. Head-on: $13.6\un{TeV}$, $120$ times more. To match the LHC, a fixed-target beam would need $E = s/2mc^2 = (13\,600)^2/(2\times0.938)\un{GeV}\approx10^{17}\un{eV}$: such energies are found only in cosmic rays.

\sol{accel}{5} LEP: $88.5\times104^4/3100\un{keV} = 3.3\un{GeV}$ per turn, $3\%$ of the beam energy, replaced every turn by the RF cavities. LHC: $88.5\times6800^4/2800\times(1/1836)^4\un{keV} = 6\un{keV}$ per turn, utterly negligible as a fraction of the energy. But the LHC has $3\times10^{14}$ protons per beam, each circulating $11\,000$ times a second: the radiated power is about $3.5\un{kW}$ per beam. Removing one watt at $1.9\un{K}$ takes about a kilowatt of refrigeration power at room temperature, so the radiation is intercepted by a ``beam screen'' at $5$--$20\un{K}$ lining the beam pipe.

\sol{accel}{6} $\sigma L = 55\times10^{-36}\un{cm^2}\times2\times10^{34}\un{cm^{-2}s^{-1}} = 1.1$ per second, or $10^7$ per year; $0.23\%$ of these is $25\,000$ diphoton events per year. But there are about a billion proton collisions per second, and ordinary processes make pairs of photons (and jets that mimic photons) at a rate thousands of times higher. The Higgs appeared as a small bump on a large, smooth background in the two-photon mass spectrum, which needed years of data, excellent energy resolution, and careful rejection of fakes. (A billion collisions a second also have to be filtered down to about a thousand recorded per second, in real time.)

\sol{accel}{7} $\cos\theta = 1/n\beta$ gives $\theta = 41^\circ$ for $\beta\approx1$. Emission requires $\beta>1/n = 0.75$, i.e.\ $\gamma>1.52$ and $\beta\gamma>1.14$: $p = \beta\gamma mc$ is $121\un{MeV}/c$ for a muon and $1.07\un{GeV}/c$ for a proton.

\sol{accel}{8} $p = 0.3\times3.8\times10 = 11.4\un{GeV}/c$. $s = L^2/8R = 1/80\un{m} = 12.5\un{mm}$. Since $s\propto1/R\propto1/p$, $\delta p/p = \delta s/s = 50\un{\mu m}/12.5\un{mm} = 0.4\%$. At higher momenta the sagitta shrinks and the resolution worsens in proportion to $p$, which is why muon spectrometers are huge and their magnets strong.

\sol{accel}{9} $\ell = 1/(n\sigma) = 1/(7\times10^{30}\times10^{-47}) = 1.4\times10^{16}\un{m}$, one and a half light-years of lead. Reines and Cowan placed a few hundred litres of water, with cadmium dissolved in it, next to the Savannah River reactor. With about $3\times10^{28}$ protons, the interaction rate $\Phi\sigma N\approx10^{17}\times10^{-47}\times3\times10^{28}$ is a few per hundred seconds. The reaction $\bar\nu_e + p\to n + e^+$ produces a distinctive pair of signals: a prompt flash from the positron annihilating into two $511\un{keV}$ photons, then a second flash a few microseconds later as the neutron is captured by cadmium. Requiring this delayed coincidence picked a few events per hour out of the background.

\sol{accel}{10} With $f_1 = f$, $f_2 = -f$: $1/F = 1/f - 1/f + d/f^2 = d/f^2$. In the other plane, $f_1 = -f$, $f_2 = f$, and again $1/F = d/f^2$. Either way, $F = f^2/d>0$. A thin lens deflects a ray in proportion to its distance from the axis. In the plane where the first lens focuses, the beam arrives at the second (defocusing) lens narrower, so it gets a weaker defocusing kick. In the other plane, the first lens defocuses, so the beam arrives wider at the focusing lens and gets a stronger focusing kick. In both planes the focusing lens acts where the beam is wide and the defocusing lens where it is narrow, so focusing wins. Alternating gradients made accelerators like the LHC possible: the beam pipe can be centimetres wide instead of metres.

\sol{accel}{11} Each beam has $E = 500\un{GeV}$. The power lost by a beam of current $I$ is $I\times\Delta E/e$, so $100\un{MW}$ at $10\un{mA}$ allows $\Delta E = 10\un{GeV}$ per turn. Then $88.5\times500^4/\rho\un{keV} = 10^7\un{keV}$ gives $\rho = 550\un{km}$: a ring over $3500\un{km}$ around, for one beam. A linear collider at $40\un{MV/m}$ needs $500\un{GeV}/(40\un{MV/m}) = 12.5\un{km}$ per beam, $25\un{km}$ in all. Because synchrotron losses grow as $E^4/\rho$, the cost of a circular electron machine grows as $E^2$ (balancing ring size against power), while a linear collider's grows only as $E$. Above a few hundred GeV, electrons must go straight.

\sol{accel}{12} If the plasma electrons, density $n$, are displaced by a distance $x$, Gauss's law gives a field $E = nex/\epsilon_0$. The wave cannot displace them by much more than its own wavelength over $2\pi$, about $c/\omega_p$, before it breaks. So $E_{\max}\sim nec/\epsilon_0\omega_p = m_ec\,\omega_p/e$, using $\omega_p^2 = ne^2/\epsilon_0m_e$. Numerically, $m_ec\,\omega_p/e = 96\sqrt{n[\mathrm{cm^{-3}}]}\un{V/m}$. For $n = 10^{18}\un{cm^{-3}}$: $10^{11}\un{V/m} = 100\un{GV/m}$, a thousand times more than RF cavities ($30$--$100\un{MV/m}$, limited by sparking at the metal walls; a plasma is already ionized and cannot break down). A TeV would take ten metres instead of ten kilometres. The hard parts are keeping the accelerated bunch small and monoenergetic, chaining stages together, and doing it efficiently, thousands of times a second.
